\documentclass[review]{jfp-epi}
\settopmatter{printfolios=true,printccs=true,printacmref=true}

%%
%% Journal information
%%
\acmJournal{PACMPL}
\acmVolume{1}
\acmNumber{OOPSLA}
\acmArticle{1}
\acmYear{2020}
\acmMonth{1}
\acmDOI{}
\startPage{1}

%%
%% Copyright information
%%
\setcopyright{none}

%%
%% Bibliography and Citation Style
%%
\bibliographystyle{ACM-Reference-Format}
\citestyle{acmauthoryear}

%%
%% Imported Packages
%%
\usepackage{booktabs}
\usepackage{subcaption}
\usepackage{amsmath}
\usepackage[all,cmtip]{xy}
\usepackage[nameinlink,noabbrev,capitalise]{cleveref}
\usepackage{csquotes}
\usepackage{booktabs}
\usepackage{microtype}
\usepackage{stmaryrd}
\usepackage{graphicx}
\usepackage{enumitem}
\usepackage{mathtools}
\usepackage{courier}
\usepackage{bussproofs}
\usepackage{tikz}
\usetikzlibrary{decorations.pathreplacing}
\usetikzlibrary{positioning}
\usetikzlibrary{calc}
\usetikzlibrary{cd}
\usetikzlibrary{math}
\usetikzlibrary{arrows}
\usetikzlibrary{patterns}
\usepackage{diagbox}
\usepackage{cmll}
\usepackage{afterpage}
\usepackage[normalem]{ulem}
\usepackage{dashbox}
\usepackage{mdframed}
\usepackage[mathscr]{euscript}
\DeclareMathAlphabet\mathbfscr{U}{eus}{b}{n}
\usepackage{colortbl}
\usepackage{anyfontsize}
\usepackage{listings}
\usepackage{alltt} % verbatim with colors
\usepackage{multicol} %for dividing into 2 columns
\usepackage{scalerel} %resize symbols
\usepackage{tipa}

%%
%% Theorem Environments
%%
\newtheorem{remark}{Remark}
\newtheorem{frameddefinition}{Haskell Report}[section]
\usepackage{mdframed} % Framed amsmath environments
\surroundwithmdframed{frameddefinition}

%%
%% Todo notes
%%
\setlength{\marginparwidth}{2cm} % This is needed, otherwise the package "todonotes" will generate a warning.
\usepackage[]{todonotes}
\newcommand{\davidb}[2][]{\todo[inline,color=red!40,author=DavidB,#1]{#2}}


%%
%% Macros
%%
\newcommand{\nonterminal}[1]{\langle \text{#1} \rangle}
\newcommand{\terminal}[1]{\textbf{#1}}

% Color definitions

\usepackage{xcolor}
%%
%% Used for Minted
%%
%% \usepackage{noto-mono}

\begin{document}

%%
%% Title information
%%
\title{The Haskell 2010 Module System}
%\subtitle{}

%%
%% Keywords
%%
\keywords{Haskell, Formal Specification}

%%
%% CCS Classification
%%

\begin{CCSXML}
  <ccs2012>
  <concept>
  <concept_id>10003752.10003753.10003754.10003733</concept_id>
  <concept_desc>Theory of computation~Lambda calculus</concept_desc>
  <concept_significance>300</concept_significance>
  </concept>
  <concept>
  <concept_id>10003752.10003790.10011740</concept_id>
  <concept_desc>Theory of computation~Type theory</concept_desc>
  <concept_significance>300</concept_significance>
  </concept>
  </ccs2012>
\end{CCSXML}

\ccsdesc[300]{Theory of computation~Lambda calculus}
\ccsdesc[300]{Theory of computation~Type theory}

%%
%% Author: David Binder
%%
\author{David Binder}
\orcid{0000-0003-1272-0972}
\email{D.Binder@kent.ac.uk}
\affiliation{
  \department{Department of Computer Science}
  \institution{University of Kent}
  \city{Canterbury}
  \country{United Kingdom}
}

%%
%% Abstract
%%
\begin{abstract}
  A formal specification of the Haskell 2010 module system.
\end{abstract}


\maketitle

%%
%% Module System
%%
\section{Module System}
\label{sec:module-system}
%%
%% Subsec: Concrete Syntax Tree
%%
\subsection{Concrete Syntax Tree}
\label{subsec:module-system:cst}

In this subsection we define the \emph{concrete syntax tree (CST)} that is the result of parsing the concrete syntax of Haskell.

%% Subsubsec: Names
\subsubsection{Names}

\[
\begin{array}{rcll}
    \nonterminal{mod-name} & \coloneqq & (\terminal{mod-name}\ s\ \overline{s}) & \text{Text}, \text{Data.Text}\\
    \nonterminal{con-name} & \coloneqq & (\terminal{con-name}\ s) & \text{Maybe}, \text{Nothing}\\
    \nonterminal{q-con-name} & \coloneqq &  (\terminal{q-con-qual}\ \nonterminal{mod-name}\ \nonterminal{con-name}) & \text{Data.Maybe.Nothing}\\
    & \mid & (\terminal{q-con-unqual}\ \nonterminal{con-name}) & \text{Nothing}\\
    \nonterminal{var-name} & \coloneqq & (\terminal{var-name}\ s) & \text{isJust}\\
    \nonterminal{q-var-name} & \coloneqq &  (\terminal{q-var-qual}\ \nonterminal{mod-name}\ \nonterminal{var-name}) & \text{Data.Maybe.isJust} \\
    & \mid & (\terminal{q-var-unqual}\ \nonterminal{var-name}) & \text{isJust} \\
    \nonterminal{op-name} & \coloneqq & (\terminal{op-name}\ s) & \text{+} \\
    \nonterminal{q-op-name} & \coloneqq &  (\terminal{q-op-qual}\ \nonterminal{mod-name}\ \nonterminal{op-name}) & \text{Prelude.(+)} \\
    & \mid & (\terminal{q-op-unqual}\ \nonterminal{op-name}) & \text{+} \\
    
\end{array}
\]
%% Subsubsec: Literals
\subsubsection{Literals}

\[
\begin{array}{rcl}
    \nonterminal{literal} &\coloneqq & (\terminal{lit-char}\ c) \\
    & \mid & (\terminal{lit-string}\ s) \\
    & \mid & (\terminal{lit-int}\ n) \\
    & \mid & (\terminal{lit-float}\ n\ n)
\end{array}
\]
%% Subsubsec: Patterns
\subsubsection{Patterns}

\[
\begin{array}{rcl}
    \nonterminal{labelled-pattern} & \coloneqq & (\terminal{label-pat}\ \nonterminal{q-var-name}\ \nonterminal{pattern}) \\
    \nonterminal{pattern} & \coloneq & (\terminal{pat-lit}\ \nonterminal{literal}) \\
    & \mid & (\terminal{pat-var}\ \nonterminal{var-name}) \\
    & \mid & (\terminal{pat-irrefutable}\ \nonterminal{pattern}) \\
    & \mid & (\terminal{pat-as}\ \nonterminal{var-name}\ \nonterminal{pattern}) \\
    & \mid & (\terminal{pat-con}\ \nonterminal{q-con-name}\ \overline{\nonterminal{pattern}}) \\
    & \mid & (\terminal{pat-record}\ \nonterminal{q-con-name}\ \overline{\nonterminal{labelled-pattern}}) \\
    & \mid & (\terminal{pat-tuple}\ \overline{\nonterminal{pattern}}) \\
    & \mid & (\terminal{pat-list}\ \overline{\nonterminal{pattern}}) \\
    & \mid & \terminal{pat-wildcard}
\end{array}
\]


%%
%% Subsubsec: Expressions
%%
\subsubsection{Expressions}

\[
\begin{array}{rcl}
    \nonterminal{qual} &\coloneq & (\terminal{qual-gen}\ \nonterminal{pattern}\ \nonterminal{expr}) \\
    &\mid & (\terminal{qual-let}\ \overline{\nonterminal{decl}}) \\
    &\mid & (\terminal{qual-guard}\ \nonterminal{expr}) \\
    \nonterminal{expr} &\coloneqq & (\terminal{exp-lit}\ \nonterminal{literal}) \\
    &\mid & (\terminal{exp-var}\ \nonterminal{q-var-name}) \\
    &\mid & (\terminal{exp-list-comp}\ \nonterminal{expr}\ \overline{\nonterminal{qual}}) \\
    &\mid & (\terminal{exp-case}\ \nonterminal{expr}\ \overline{\nonterminal{alt}}) \\
    &\mid & (\terminal{exp-if-then-else}\ \nonterminal{expr}\ \nonterminal{expr}\ \nonterminal{expr}) \\
    &\mid & (\terminal{exp-tuple}\ \overline{\nonterminal{expr}}) \\
    &\mid & (\terminal{exp-list}\ \overline{\nonterminal{expr}}) \\
    &\mid & (\terminal{exp-record}\ \nonterminal{q-con-name}\ \ldots) \\
    &\mid & (\terminal{exp-record-update}\ \nonterminal{expr}\ \ldots) \\
\end{array}
\]
%% Subsubsec: Declarations
\subsubsection{Declarations}
\[
\begin{array}{rcl}
    \nonterminal{fixity} &\coloneqq & \terminal{infixl} \mid \terminal{infixr} \mid \terminal{infix} \\
    \nonterminal{local-decl} &\coloneqq & (\terminal{local-decl-fixity}\ \nonterminal{fixity}\ n\ \overline{\nonterminal{...}}) \\
    & \mid & (\terminal{local-decl-sig}\ \overline{\nonterminal{var-name}}\ \nonterminal{type}) \\
    & \mid & (\terminal{local-decl-eq}  \ldots) \\
    \nonterminal{top-decl} & \coloneqq & (\terminal{top-decl-data} \ldots) \\
    & \mid & (\terminal{top-decl-newtype} \ldots) \\
    & \mid & (\terminal{top-decl-type} \ldots) \\
    & \mid & (\terminal{top-decl-default} \ldots) \\
    & \mid & (\terminal{top-decl-class} \ldots) \\
    & \mid & (\terminal{top-decl-instance} \ldots) \\
    & \mid & (\terminal{top-decl-foreign} \ldots) \\
    & \mid & (\terminal{top-decl-local}\ \nonterminal{local-decl})
\end{array}
\]
%% Subsubsec: Modules
\subsubsection{Modules}

\[
\begin{array}{rcl}
    \nonterminal{export} &\coloneqq & \ldots \\
    \nonterminal{exports} &\coloneqq & \terminal{none} \mid (\terminal{some}\ \overline{\nonterminal{export}})\\[0.2cm]
    \nonterminal{import-item} & \coloneq & \ldots \\
    \nonterminal{hiding} & \coloneqq & \terminal{hiding} \mid \terminal{exposed}\\
    \nonterminal{qualified} & \coloneqq & \terminal{qualified} \mid \terminal{unqualified}\\
    \nonterminal{import-as} & \coloneq & \terminal{none} \mid (\terminal{some}\ \nonterminal{modname}) \\
    \nonterminal{import} & \coloneqq & (\terminal{import}\ \nonterminal{mod-name}\ \nonterminal{hiding}\ \nonterminal{qualified}\ \nonterminal{import-as}\ \overline{\nonterminal{import-item}}) \\[0.2cm]
    \nonterminal{module} &\coloneqq & (\terminal{module}\ \nonterminal{mod-name}\ \nonterminal{exports}\ \overline{\nonterminal{import}}\ \overline{\nonterminal{top-decl}})
\end{array}
\]

%%
%% Subsec: Renamed Syntax Tree
%%
\subsection{Renamed Syntax Tree}
\label{subsec:module-system:rst}

In this subsection we define the \emph{renamed syntax tree (RST)} that is the result of name resolution.

\subsubsection{Modules}

\[
\begin{array}{rcl}
    \nonterminal{module} &\coloneqq & (\terminal{module}\ \overline{\nonterminal{topdecl}})
\end{array}
\]


%%
%% Related Work
%%
\section{Related Work}
\label{sec:related-work}
\begin{itemize}
    \item \cite{Faxen2002staticsemantics}
    \item \cite{Diatchki2002haskellmodules}
\end{itemize}

%%
%% Bibliography
%%
\bibliography{bibliography/bibliography.bib}

\end{document}
